% Part: first-order-logic
% Chapter: lindstrom
% Section: ls-property

\documentclass[../../../include/open-logic-section]{subfiles}

\begin{document}

\olfileid{mod}{lin}{lsp}

\olsection{محدود گواہی تے L\"owenheim--Skolem دیاں خاصیتاں}

ہن اسیں تجریدی منطقاں لئی محدود گواہی تے
L\"owenheim--Skolem دیاں سیدھیاں توسیعاں دیندے آں۔

\begin{defn}
اک تجریدی منطق $\tuple{L, \models_L}$ وچ \emph{محدود گواہی
  دی خاصیت} اے جے ہر سیٹ $\Gamma$، جو
$L(\Lang{L})$ دے !!{sentence}s دا ہووے، اوہدوں قابلِ اطمینان ہووے جدوں اوہدا
ہر محدود $\Gamma_0 \subseteq \Gamma$ قابلِ اطمینان ہووے۔
\end{defn}

\begin{defn}
$\tuple{L, \models_L}$ وچ \emph{نزولی L\"owenheim--Skolem
  دی خاصیت} اے جے ہر قابلِ اطمینان $\Gamma$ دا
!!a{enumerable} نمونہ موجود ہووے، یعنی محدود یا گنتی وار
لامتناہی نمونہ۔
\end{defn}


\olref[bas][pis]{defn:partialisom} وچ جزوی ہم شکلی دا تصور خالص
«الجبری» اے، یعنی اوہ زبان دے !!{sentence}s نوں
نہیں، بلکہ صرف !!{structure}s دی !!{language}~$\Lang{L}$
وچ دتیاں غیر منطقی علامتاں نوں ویکھ کے بیان ہوندا اے؛ ایس لئی ایہ
تجریدی منطقاں اُتے وی لاگو ہوندا اے۔ پہلے درجے دی منطق لئی
\olref[bas][pis]{thm:p-isom2} توں سانوں پتہ اے کہ جے دو
!!{structure}s جزوی طور اُتے ہم شکل ہون تاں اوہ جملہ وار
ہم ارز نیں۔ اوہ ثبوت تجریدی منطقاں تک نہیں پہنچدا، کیوں جے
من مانے $!E \in L(\Lang{L})$ لئی !!{formula}s اُتے
استقرا دستیاب ہونا ضروری نہیں؛ پر جے L\"owenheim--Skolem دی
خاصیت ہووے، تاں قضیہ پھر وی درست اے۔

\begin{thm}
\ollabel{thm:abstract-p-isom}
فرض کرو $\tuple{L, \models_L}$ اک معمول مطابق منطق اے جس وچ
L\"owenheim--Skolem دی خاصیت اے۔ تاں ہر دو !!{structure}s
جو جزوی طور اُتے ہم شکل ہون،
$\tuple{L,
  \models_L}$ وچ جملہ وار ہم ارز نیں۔
\end{thm}

\begin{proof}
فرض کرو $\Struct{M} \simeq_p \Struct{N}$، پر کسے $!E$ لئی
$\Struct{M} \models_L !E$ جدوں کہ
$\Struct{N} \not\models_L !E$۔ ہم شکلی دی خاصیت توں اسیں
ہم شکل نقلاں ورت سکدے آں۔ پر مستقل نشانیاں ہون تاں
$\Domain{M}$ تے $\Domain{N}$ نوں پوری طرح وکھرا نہیں رکھنا:
پرانیاں مستقل نشانیاں دیاں متناظر قدراں نوں اکّو عنصر بنا کے پہچانو،
تے باقی عناصر نوں وکھرے نشان دیو۔ اوہ سانجھا سیٹ $K$ آکھو۔
اگے-پیچھے خاصیت نال ہر محدود لسٹ دیاں نامی قدراں کسے جزوی نقشے
وچ آ سکدیاں نیں؛ سو دوہاں ساختاں دے مستقل برابر ہون دے نمونے
اکّو نیں، تے سانجھیاں نامی قدراں اُتے ہر محمول دی تعبیر وی اکّو اے۔
ایس لئی ایہ پہچان اک درست سانجھا حصہ بناندی اے۔
اسیں
$!E \in
L(\Lang{L})$ لئی اصل !!{language}~$\Lang{L}$ برقرار رکھاں گے؛
توسیع دی خاصیت صرف محدود علامتاں اُتے صدق دی وابستگی دیندی اے،
ایس توں جملے دی محدود زبان وچ آمد اپنے آپ نہیں نکلدی۔ ہِٹھاں
پہلے درجے دے جملیاں دا سیٹ ورتاں گے، جس لئی نزولی خاصیت کافی اے۔
$\mathcal{I}$، $\Struct{M}$ تے $\Struct{N}$ وچکار جزوی
ہم شکلیاں دا اک غیر خالی، اگے-پیچھے خاصیت والا سیٹ ہووے۔
عمومیت گنوائے بغیر ایہ وی فرض
کرو کہ جے $p
\in \mathcal{I}$ تے $q \subseteq p$، تاں
$q \in \mathcal{I}$ وی۔ ہر ایسے نقشے نوں تمام سانجھیاں نامی
قدراں اُتے شناختی نقشے نال ملا کے $j_p=p\cup\operatorname{id}_K$
لکھو۔ کسے نامی قدر نوں اصل نقشے وچ شامل کرن دی توسیع اوہنوں
اوہدی متناظر قدر اُتے ای بھیجدی اے؛ یک بہ یک ہون کارن کوئی ہور
عنصر اوس قدر اُتے نہیں جا سکدا۔ سو ایہ ملاپ وی جزوی یک بہ یک
نقشہ اے تے سارے محمول محفوظ رکھدا اے:
کسے محمول نوں پرکھن لئی صرف محدود ٹپل دی لوڑ اے، تے اوس ٹپل
دیاں نامی قدراں نوں اصل اگے-پیچھے سیٹ دے نقشے وچ شامل کیتا جا سکدا اے۔
نویں نقشے محدود ہون ضروری نہیں؛ پر اوہناں دا غیر خالی خاندان ہن وی
دونوں سمتاں وچ ہر نواں عنصر شامل کر سکدا اے۔

$\Domain{M}^{<\omega}$، $\Domain{M}$ دے !!{element}s دیاں
محدود ترتیب وار لڑیّاں دا سیٹ اے؛ ساخت بناندیاں لڑیّاں دیاں
وکھرے نشان لائیاں نقلاں ورتو، تاکہ عنصر تے لڑی دے کوڈ آپس وچ
نہ رلن، تے دوہاں پاسیاں دیاں خالی لڑیّاں دے کوڈ وی وکھرے ہون۔ $S$،
$\Domain{M}^{<\omega}$ اُتے تین مقامی رشتہ ہووے جو لڑیّاں
جوڑن نوں ظاہر کردا اے؛ یعنی جے $\mathbf{a}, \mathbf{b},
\mathbf{c} \in \Domain{M}^{<\omega}$، تاں
$S(\mathbf{a}, \mathbf{b},
\mathbf{c})$ اوہدوں تے صرف اوہدوں درست اے جدوں
$\mathbf{c}$، $\mathbf{a}$ تے $\mathbf{b}$ نوں جوڑ
کے بنی ہووے۔ تے $T$ اوہ تین مقامی رشتہ ہووے جس لئی
$T(\mathbf{a}, b, \mathbf{c})$، $b \in M$ تے
$\mathbf{a}, \mathbf{c} \in \Domain{M}^{<\omega}$ لئی
اوہدوں تے صرف اوہدوں درست اے جدوں $\mathbf{a} =
a_1, \dots a_n$ تے $\mathbf{c} = a_1, \dots a_n, b$۔
نویں تین مقامی !!{predicate}s $P$ تے $Q$ چنو، تے
!!{structure} $\Struct{M}^*$ بناؤ جس دا دائرہ
$\Domain{M} \cup \Domain{M}^{<\omega}$ اے،
جس وچ $\Struct{M}$ ذیلی ساخت اے، تے جس وچ $P$ تے $Q$
دی تعبیر بالترتیب جوڑن دے رشتے $S$ تے $T$ نیں
(اس لئی $\Struct{M}^*$،
!!{language} $\Lang{L} \cup \{ P, Q\}$ وچ اے)۔

$\Domain{N}^{<\omega}$، $S'$، $T'$، $P'$، $Q'$ تے
$\Struct{N}^*$ اوہو جیہے طریقے نال مقرر کرو۔ چونکہ مفروضے
مطابق $\Struct{M} \simeq_p
\Struct{N}$، ایس لئی اک رشتہ $I$،
$\Domain{M}^{<\omega}$ تے
$\Domain{N}^{<\omega}$ وچکار اے جس لئی
$I(\mathbf{a}, \mathbf{b})$ اوہدوں تے صرف اوہدوں درست اے
جدوں برابر لمبائی دیاں لڑیّاں $\mathbf{a}$ تے $\mathbf{b}$
دے متناظر عناصر نوں کسے نقشے $p\in\mathcal{I}$ نال ملایا جا سکدا ہووے؛
\olref[bas][pis]{defn:partialisom} دی اگے-پیچھے شرط
پورے خاندان دی شرط اے، اکیلے ٹپل دی نہیں۔ ہن $\Struct{C}$
نوں اوہ !!{structure} آکھو جس دا
!!{domain}، $\Struct{M}^*$ تے $\Struct{N}^*$ دے
!!{domain}s دا اتحاد، نال ہر نقشے $j_p$ دا اک وکھرا کوڈ، اے۔
اصل زبان دے محمول دوہاں پاسیاں اُتے اصل تعبیر رکھن، تے باقی
مخلوط ٹپلاں اُتے جھوٹے رکھے جا سکدے نیں؛ سانجھا حصہ $K$ اُتے
تعبیراں پہلے ای رلدیاں نیں۔ اپنے اپنے متعلقہ زبان والے گھٹائے روپ
وچ $\Struct{M}^*$ تے $\Struct{N}^*$ ذیلی !!{structure}s نیں۔
ایس ساخت دی
!!{language} وچ اک وادھو دو مقامی !!{predicate}~$R$ اے
جس دی تعبیر رشتہ~$I$ اے، نالے اوہ !!{predicate}s
جو !!{domain}s $\Domain{M}^*$ تے~$\Domain{N}^*$
نوں ظاہر کردے نیں۔ ہور نویں یک مقامی محمول $U$ تے $V$
صرف اصل دو دائرے چنّن، یک مقامی محمول $H$ نقشیاں دے کوڈ چنّے،
تے تین مقامی محمول $G$ لئی $G(h,a,b)$ دا مطلب ہووے کہ
کوڈ $h$ والا نقشہ عنصر $a$ نوں $b$ نال ملاندا اے۔

\begin{figure}[h]
  \centering
  \begin{tikzpicture}[node distance=2cm, auto, thick, >=stealth']
    \draw [rounded corners] (0,0) -- (8,0) -- (8,4) -- (0,4) --  cycle;
    \draw (2,2) circle (0.5cm);
    \draw (2,2) circle (1.25cm);
    \draw (6,2) circle (0.5cm);
    \draw (6,2) circle (1.25cm);
    \path node at (0.75,3.5) {\large $\Struct{C}$};
    \path node at (2,2) {\large $\Struct{M}$};
    \path node at (6,2) {\large $\Struct{N}$};
    \path node at (3.5,1) {\large $\Struct{M}^*$};
    \path node at (7.5,1) {\large $\Struct{N}^*$};
    \node (Idom) at (2.8,2) {};
    \node (Irng) at (5.2,2) {};
    \draw[<->, bend left] (Idom) to node {\large $I$} (Irng) ;
  \end{tikzpicture} 
  \caption{!!{structure}~$\Struct{C}$ تے اوہدے اندرلی
    جزوی ہم شکلی دا خاکہ۔ وکھرے دائرے مستقل نشانیاں توں خالی حال
    دکھاندے نیں؛ مستقل ہون تاں دوہاں پاسیاں دیاں سانجھیاں نامی قدراں
    نوں اکّو عنصر سمجھنا اے۔ نقشیاں دے وادھو کوڈ بیرلے دائرے وچ نیں،
    پر ایس خاکے وچ الگ نہیں دکھائے گئے۔}
\end{figure}

بنیادی مشاہدہ ایہ اے کہ !!{structure}~$\Struct{C}$ دی
!!{language} وچ اک \emph{تجریدی} !!{sentence} $!D_1$
موجود اے، جو $\Struct{C}$ وچ سچ اے تے آکھدا اے کہ
$\Struct{M} \models_L !E$ تے
$\Struct{N} \not\models_L !E$ (ایس لئی نسبتی بنانے
دی خاصیت تے بولین بندش درکار نیں)۔ ایہ جملہ پہلے نسبتی روپ
توں بنے گا: پہلا روپ $U$ تک، تے نفی والا دوجا روپ $V$ تک۔
دونوں چنے دائرے غیر خالی نیں تے ہر پرانے مستقل دی قدر دوہاں وچ اے۔
توسیع دی خاصیت وچ صدق صرف اوہ محدود گھٹایا روپ ویکھدا اے؛
لہٰذا جملہ جس زبان وچ دستیاب ہووے، اوہدے ہور نویں علامتاں نال
ودھائے روپ اُتے وی اوہو صدق رہندا اے۔ ایہ صدقی بات اے، جملے نوں
محدود زبان دا قرار دین دی بات نہیں۔ من مانے تجریدی اصل جملے لئی
ایس بنے جملے نوں پہلے درجے دا نہیں کہہ سکدے۔
اصل دے واحد \emph{پہلے درجے دے} !!{sentence} $!D_2$ دی تھاں
پہلے درجے دے جملیاں دا سیٹ $\Delta_2$ ورتو۔ ایس وچ چنے دائرے
تے کوڈ غیر خالی ہون، نقشے جزوی یک بہ یک ہون، تے اگے-پیچھے
توسیع لئی ایہ جملے شامل نیں:
\[
\begin{gathered}
\exists a\,U(a),\qquad \exists b\,V(b),\qquad \exists h\,H(h),\\
\forall h,a,b\,(G(h,a,b)\to(H(h)\land U(a)\land V(b))),\\
\forall h,a,b,b'\,((G(h,a,b)\land G(h,a,b'))\to b=b'),\\
\forall h,a,a',b\,((G(h,a,b)\land G(h,a',b))\to a=a'),\\
\forall h,a\,((H(h)\land U(a))\to
\exists k,b\,(H(k)\land V(b)\land G(k,a,b)\land
\forall u,v\,(G(h,u,v)\to G(k,u,v)))),\\
\forall h,b\,((H(h)\land V(b))\to
\exists k,a\,(H(k)\land U(a)\land G(k,a,b)\land
\forall u,v\,(G(h,u,v)\to G(k,u,v)))).
\end{gathered}
\]
ہر پرانے مستقل $c$ لئی جملے $U(c)\land V(c)$ تے
$\forall h\,(H(h)\to G(h,c,c))$ شامل کرو۔ ہر پرانے
$n$-مقامی محمول $A$ لئی ایہ جملہ شامل کرو:
\[
\forall h,a_1,\ldots,a_n,b_1,\ldots,b_n\,
\bigl((H(h)\land\bigwedge_{i=1}^{n}G(h,a_i,b_i))\to
(A(a_1,\ldots,a_n)\leftrightarrow A(b_1,\ldots,b_n))\bigr).
\]
ایہ صرف پہلے درجے دے جملیاں دا سیٹ اے، زبان محدود ہون دی شرط
نہیں؛ شناختی رشتے دی حفاظت نقشے دی یک بہ یک ہون دی شرط توں آندی اے۔
عام پہلے درجے دے ایہ جملے معمول مطابق تجریدی منطق وچ دستیاب نیں۔
ایہ سارا سیٹ $\Struct{C}$ وچ سچ اے: ہر کوڈ اصل $j_p$ دا اے،
توسیع اصل اگے-پیچھے خاندان دی توسیع توں ملدی اے۔
L\"owenheim--Skolem دی خاصیت جملیاں دے پورے قابلِ اطمینان سیٹ
$\{!D_1\}\cup\Delta_2$ نوں !!a{enumerable}
نمونہ $\Struct{C}_0$ دیندی اے۔ اوہدے $U$ تے $V$ والے دائرے
چن کے اصل زبان دیاں ذیلی ساختاں $\Struct{M}_0$ تے
$\Struct{N}_0$ لوو؛ اوہ غیر خالی تے قابلِ شمار نیں، تے
$\Struct{M}_0 \models_L !E$ تے
$\Struct{N}_0
\not\models_L !E$۔ کوڈاں دے نقشیاں دیاں ساریاں محدود پابندیاں
اک غیر خالی جزوی ہم شکلی خاندان بناندیاں نیں: کسے محدود نقشے
دا اصل کوڈ چن کے اُتّے لکھی توسیع نال نواں عنصر شامل کرو، تے
فیر صرف لوڑ جوگی محدود پابندی لوو۔ محمولاں تے مستقلاں دی حفاظت
درج جملیاں توں آندی اے۔ ایس طرح $\Struct{M}_0 \simeq_p \Struct{N}_0$
دا اصل دلیل والا اگے-پیچھے طریقہ ہن نئی دو ساختاں لئی وی حاصل اے؛
رشتہ~$I$ دی اصل کوڈ بندی نوں نقشیاں دے گراف نال صریح کیتا گیا اے۔
پر !!{enumerable}، جزوی طور اُتے
ہم شکل !!{structure}s دراصل
\olref[bas][pis]{thm:p-isom1} توں ہم شکل ہوندیاں نیں؛
ایہ معمول مطابق تجریدی منطقاں دی ہم شکلی دی خاصیت دے خلاف اے۔
\end{proof}

\end{document}
